% !TeX root = ../../tesis.tex
%%%%%%%%%%%%%%%%%%%%%%%%%%%%%%%%%%%%%%%%%%%%%%%%%%%%%%%%%%%%%%%%%%%
% §6.3.5 — CENTRALIDAD DE INTERMEDIACIÓN B(v) EN LOS CASOS CON ANILLO
%
% Datos (corrida 2026-10-04; bv_comparativo.csv):
%   Tacubaya 0.2232 → 0.2163 (Caso 2, −3.1 %) → 0.1569 (Caso 3, −29.7 %)
%   Caso 2: ninguna estación del anillo en el top 50; MB conserva los
%     10 nodos de la red actual (Periférico Oriente/Tláhuac 0.061, #69)
%   Caso 3: 9/10 y 40/50 del top son del anillo; top 5 = Tacubaya,
%     Periférico Oriente/Tláhuac 0.1550, Rómulo O'Farril 0.1310,
%     Barranca del Muerto/Periférico 0.1284, Constitución de
%     1917/Periférico 0.1274; 208 nodos pierden centralidad.
%   Gini 0.8007 → 0.8012 → 0.8134 (0.199 · 0.199 · 0.187, Crítico).
% Criterio de las figuras (decisión C6-5):
%   comparativos = orden por nodo; Dashboard_Betweenness_* = 10
%   nombres con mayor suma de B(v) (como tab:bv-top10-agregado).
% Estructura: entrada · Caso 2 · Caso 3 · lectura comparativa.
% Figuras candidatas:
%   .../Clasificacion_Garibelt_Betweenness_Centrality_MB / _Metro
%   Tableau_Resources/RedConAnillo/Escenario_MB|Metro/Dashboard_Betweenness_*
%   Qgis_Resources/Comparativo_CG/Comparativo_Garibelt_Betweenness
%   Tableau_Resources/Comparativo_CG/Dashboard_Comparativo_Intermediacion_Bv
%%%%%%%%%%%%%%%%%%%%%%%%%%%%%%%%%%%%%%%%%%%%%%%%%%%%%%%%%%%%%%%%%%%
